\chapter{KV Cache 机制}
\label{chap:kv-cache}

\section{自回归推理流程}
\label{sec:autoregressive-inference-flow}

前面几篇主要围绕训练展开：从 Transformer 结构到分布式训练，再到 3D 并行和 ZeRO。现在开始进入第四篇：\textbf{推理架构优化 Serving Infra}。训练系统关心的是如何在大量 GPU 上稳定、高吞吐地完成参数更新；推理系统关心的是如何在在线请求中以低延迟、低成本、高并发生成 token。

大语言模型推理和普通深度学习推理有本质差异。一个图像分类模型通常一次 forward 得到一个分类结果；一个推荐模型通常一次 forward 得到候选排序；而自回归语言模型必须不断生成下一个 token，并把这个 token 追加到上下文，再生成下一个 token。也就是说，LLM 推理不是一次计算，而是一串强依赖的迭代过程：

\[
x_{t+1}\sim p_{\theta}(x_{t+1}\mid x_{\le t}),
\]
\[
x_{t+2}\sim p_{\theta}(x_{t+2}\mid x_{\le t+1}),
\]
\[
\cdots
\]

对于 Decoder-only LLM，每一步生成都依赖前面已经生成的全部 token。这个依赖使得推理阶段天然分成两个部分：

\begin{itemize}
\item \textbf{Prefill 阶段}：处理用户输入 prompt，计算 prompt 中所有 token 的 hidden states，并建立 KV Cache；
\item \textbf{Decode 阶段}：每次只处理一个新 token，读取历史 KV Cache，生成下一个 token，并追加新的 K/V。
\end{itemize}

KV Cache 正是 LLM 推理性能优化的核心。没有 KV Cache，decode 每一步都要重复计算历史 token 的 Key 和 Value；有了 KV Cache，系统只需要为新 token 计算新的 K/V，同时复用历史 K/V。这个优化看似简单，却决定了 LLM serving 的显存结构、调度策略、batch 管理、长上下文能力和吞吐上限。

\subsection{prefill 阶段}
\label{subsec:prefill-stage}

Prefill 阶段处理的是用户输入 prompt。假设 prompt token 序列为：

\[
x_1,x_2,\ldots,x_S.
\]

模型需要一次性计算这些 token 的 hidden states，并得到最后一个位置的 logits，用于采样第一个输出 token：

\[
p(x_{S+1}\mid x_1,\ldots,x_S).
\]

在 prefill 中，输入序列长度为 $S$，模型可以像训练 forward 一样并行处理整个 prompt。每一层 attention 需要计算：

\[
\mat{Q}\in\mathbb{R}^{B\times n_q\times S\times d_h},
\]
\[
\mat{K}\in\mathbb{R}^{B\times n_{kv}\times S\times d_h},
\]
\[
\mat{V}\in\mathbb{R}^{B\times n_{kv}\times S\times d_h}.
\]

对于标准 causal attention，prefill 阶段每个 token 只能关注它之前的位置，因此仍需要 causal mask。但与训练一样，所有位置可以在同一次 forward 中并行计算。Attention 的 score matrix 大小约为：

\[
B\times n_q\times S\times S.
\]

因此 prefill 的计算复杂度包含 $S^2$ 项。序列越长，prefill 越昂贵。尤其当 prompt 很长时，TTFT，即 time to first token，往往主要由 prefill 决定。

Prefill 阶段完成后，系统会把每一层的 Key 和 Value 保存下来：

\[
\mat{K}^{(\ell)}*{1:S},
\quad
\mat{V}^{(\ell)}*{1:S},
\quad
\ell=0,\ldots,L-1.
\]

这些保存下来的张量就是 KV Cache。它们会在后续 decode 中被反复读取。

从系统角度看，prefill 阶段具有以下特点：

\begin{itemize}
\item \textbf{输入长度可能很长}：prompt 可以是几百、几千甚至更多 token；
\item \textbf{计算密度较高}：QKV projection、FFN 和 attention 都能形成较大 GEMM；
\item \textbf{更接近训练 forward}：一次处理完整序列；
\item \textbf{会建立 KV Cache}：为每层、每个请求保存历史 K/V；
\item \textbf{决定首 token 延迟}：用户体感中的“模型多久开始回答”主要受 prefill 影响。
\end{itemize}

\subsection{decode 阶段}
\label{subsec:decode-stage}

Decode 阶段从模型生成第一个输出 token 后开始。假设已经有上下文：

\[
x_1,x_2,\ldots,x_t.
\]

下一步只需要处理最新 token $x_t$，生成：

\[
x_{t+1}.
\]

在第 $\ell$ 层，对最新 token 计算 query、key、value：

\[
\vect{q}_t^{(\ell)},
\quad
\vect{k}_t^{(\ell)},
\quad
\vect{v}_t^{(\ell)}.
\]

然后把新的 $\vect{k}_t,\vect{v}_t$ 追加到该层的 KV Cache 中：

\[
\mat{K}_{cache}^{(\ell)}[t]
\vect{k}*t^{(\ell)},
\]
\[
\mat{V}*{cache}^{(\ell)}[t]
\vect{v}_t^{(\ell)}.
\]

Attention 输出为：

\[
\vect{o}_t^{(\ell)}
\softmax
\left(
\frac{
\vect{q}*t^{(\ell)}
\left(\mat{K}*{cache}^{(\ell)}[1:t]\right)^{\mathsf{T}}
}
{\sqrt{d_h}}
\right)
\mat{V}_{cache}^{(\ell)}[1:t].
\]

这里非常关键：decode 每一步只为当前 token 计算 query，但需要读取历史所有 Key 和 Value。因此 decode 的每一步计算量随当前上下文长度 $t$ 线性增长，而不是像 prefill 那样一次计算 $S\times S$ 的矩阵。

Decode 阶段的系统特征与 prefill 不同：

\begin{itemize}
\item \textbf{每步只处理一个新 token}；
\item \textbf{强自回归依赖}：下一步必须等待当前 token 采样完成；
\item \textbf{GEMM 尺寸小}：batch 内每个请求通常只增加一个 token；
\item \textbf{频繁读取 KV Cache}：每一步、每一层都读取历史 K/V；
\item \textbf{通常 memory-bound}：吞吐受 HBM 带宽和 cache layout 影响很大；
\item \textbf{决定持续生成速度}：用户体感中的 tokens per second 主要受 decode 影响。
\end{itemize}

这也是 LLM serving 的难点：prefill 像一个大 batch forward，decode 像大量小步迭代；两者在算子形状、显存访问、调度目标上完全不同。

\subsection{token-by-token generation}
\label{subsec:token-by-token-generation}

完整生成流程可以写成如下循环：

\[
\text{prompt}
\xrightarrow{\text{prefill}}
\text{first logits}
\xrightarrow{\text{sample}}
x_{S+1}
\xrightarrow{\text{decode}}
x_{S+2}
\xrightarrow{\text{decode}}
\cdots
\]

每次 decode 后，模型得到下一个 token 的 logits：

\[
\vect{z}_t\in\mathbb{R}^{V},
\]
其中 $V$ 是词表大小。随后经过采样策略得到 token：

\[
x_{t+1}
\operatorname{Sample}
(
\vect{z}_t;
T,\ top\text{-}k,\ top\text{-}p,\ \ldots
).
\]

生成循环通常会在以下条件之一满足时结束：

\begin{itemize}
\item 生成了 EOS token；
\item 命中 stop words 或 stop sequence；
\item 达到最大生成长度；
\item 请求被用户取消；
\item server 触发超时或资源限制；
\item 安全策略或工具调用逻辑中断生成。
\end{itemize}

下面给出一个极简自回归推理伪代码。它强调 prefill、decode 和 KV Cache 的关系，不代表高性能实现。

\begin{lstlisting}[language=Python,caption={自回归推理中 prefill、decode 与 KV Cache 的极简伪代码},label={lst:autoregressive-generation-kv-cache}]
def generate(model, input_ids, max_new_tokens, sampler):
"""
input_ids: [batch, prompt_len]
"""
# Prefill: process the whole prompt and build KV cache.
logits, kv_cache = model.forward_prefill(input_ids)

next_token = sampler(logits[:, -1, :])
output_tokens = [next_token]

# Decode: one token per iteration.
for _ in range(max_new_tokens - 1):
    logits, kv_cache = model.forward_decode(
        input_ids=next_token[:, None],
        kv_cache=kv_cache,
    )

    next_token = sampler(logits[:, -1, :])
    output_tokens.append(next_token)

    if should_stop(next_token):
        break

return concat_tokens(output_tokens)

\end{lstlisting}

真实 serving 系统会比这复杂得多。它需要把多个请求放到同一个 batch 中，支持请求动态加入和退出，管理不同长度的 KV Cache，处理 prefix cache、分页分配、并行采样、流式输出、超时、取消、优先级、SLA 和多 GPU 通信。

\subsection{为什么推理不同于训练}
\label{subsec:why-inference-differs-from-training}

训练与推理都执行 Transformer forward，但系统形态差异很大。训练通常有固定 batch、固定 sequence length、完整 forward+backward、规则的张量形状和长时间稳定运行。推理则面向在线请求，输入长度和输出长度高度动态，且 decode 阶段存在严格 token 依赖。

训练中，一次 step 可能处理：

\[
B\times S
\]
个 token，并对所有 token 计算 loss。推理 prefill 处理 prompt，而 decode 每步通常只处理每个请求的一个 token。

训练主要瓶颈包括：

\begin{itemize}
\item 大 GEMM 计算吞吐；
\item 激活显存；
\item backward 与 optimizer；
\item DP/TP/PP 通信；
\item checkpoint 和数据加载。
\end{itemize}

推理主要瓶颈包括：

\begin{itemize}
\item KV Cache 显存；
\item decode 阶段 HBM 带宽；
\item 请求调度；
\item batch 内长度不均；
\item TTFT 与 TPOT；
\item cache 碎片；
\item 多租户和 SLA。
\end{itemize}

\begin{table}[H]
\centering
\small
\caption{LLM 训练与推理的系统差异}
\label{tab:training-vs-inference}
\begin{tabular}{p{0.22\textwidth}p{0.34\textwidth}p{0.34\textwidth}}
\toprule
\textbf{维度} & \textbf{训练} & \textbf{推理} \
\midrule
输入形态
& 大 batch、固定或近似固定序列长度
& 请求长度动态，prompt 和 output 长度差异大 \
\midrule
计算流程
& forward + backward + optimizer step
& prefill + 多轮 decode \
\midrule
并行性
& 一个 step 内大量 token 可并行
& decode token-by-token，强依赖上一 token \
\midrule
主要显存
& 参数、梯度、优化器状态、激活
& 参数、KV Cache、临时 attention buffer \
\midrule
主要瓶颈
& GEMM、激活、分布式通信
& HBM 带宽、KV Cache 管理、调度 \
\midrule
关键指标
& tokens/s、MFU、loss、收敛稳定性
& TTFT、TPOT、吞吐、尾延迟、成本/token \
\bottomrule
\end{tabular}
\end{table}

\begin{figure}[H]
\centering
\begin{tikzpicture}[
font=\small,
box/.style={draw, rounded corners=3pt, minimum width=1.8cm, minimum height=0.7cm, align=center},
token/.style={draw, rounded corners=2pt, minimum width=0.55cm, minimum height=0.4cm, align=center},
arrow/.style={-{Latex[length=2.2mm]}, thick}
]

\node[font=\bfseries] at (0,3.4) {Prefill 与 Decode 阶段对比};

% Prefill
\node[font=\bfseries] at (-3.8,2.55) {Prefill};
\foreach \i/\x in {1/-5.2,2/-4.55,3/-3.9,4/-3.25,5/-2.6} {
  \node[token, fill=blue!10] (p\i) at (\x,1.8) {$x_{\i}$};
}
\node[box, fill=orange!14] (pf) at (-3.9,0.75) {Full prompt\\forward};
\node[box, fill=green!10] (pkv) at (-3.9,-0.25) {Build KV Cache};
\node[box, fill=purple!10] (plogit) at (-3.9,-1.25) {First logits};

\foreach \i in {1,2,3,4,5} {
  \draw[arrow] (p\i) -- (pf);
}
\draw[arrow] (pf) -- (pkv);
\draw[arrow] (pkv) -- (plogit);

% Decode
\node[font=\bfseries] at (3.8,2.55) {Decode};
\node[token, fill=blue!10] (dnew) at (2.25,1.8) {$x_t$};
\node[box, fill=orange!14] (df) at (3.8,1.25) {One-token\\forward};
\node[box, fill=green!10] (dkv) at (5.35,1.8) {Read old KV\\Append new KV};
\node[box, fill=purple!10] (dlogit) at (3.8,0.25) {Next logits};
\node[token, fill=yellow!18] (sample) at (3.8,-0.75) {$x_{t+1}$};
\node[box, fill=red!8] (loop) at (3.8,-1.6) {repeat};

\draw[arrow] (dnew) -- (df);
\draw[arrow] (dkv) -- (df);
\draw[arrow] (df) -- (dlogit);
\draw[arrow] (dlogit) -- (sample);
\draw[arrow] (sample) -- (loop);
\draw[arrow] (loop.west) .. controls (1.7,-1.1) and (1.7,1.6) .. (dnew.west);

\draw[dashed, thick] (0,-2.0) -- (0,2.8);

\node[align=left, text width=10.6cm] at (0,-2.75) {
  Prefill 一次处理完整 prompt，计算密度高；Decode 每次只处理一个新 token，
  但每层都要读取不断增长的 KV Cache，因此更容易受显存带宽和 cache 管理影响。
};

\end{tikzpicture}
\caption{Prefill 与 Decode 阶段的计算流程对比}
\label{fig:prefill-vs-decode}
\end{figure}

\section{KV Cache 的基本原理}
\label{sec:kv-cache-basics}

KV Cache 的核心问题是：为什么 Key 和 Value 可以缓存？为什么 Query 不缓存？缓存的 shape 是什么？显存如何估算？这些问题决定了后续所有推理优化。

\subsection{Key / Value 为什么可以缓存}
\label{subsec:why-cache-key-value}

考虑 Decoder-only Transformer 的第 $\ell$ 层。对某个 token 的 hidden state $\vect{h}_i^{(\ell)}$，线性投影得到：

\[
\vect{q}_i^{(\ell)}=\vect{h}_i^{(\ell)}\mat{W}_Q^{(\ell)},
\]
\[
\vect{k}_i^{(\ell)}=\vect{h}_i^{(\ell)}\mat{W}_K^{(\ell)},
\]
\[
\vect{v}_i^{(\ell)}=\vect{h}_i^{(\ell)}\mat{W}_V^{(\ell)}.
\]

在自回归 decode 中，当已经生成了 token $x_1,\ldots,x_t$ 后，未来生成 $x_{t+1}$ 不会改变过去 token 的 hidden states。原因是 causal mask 保证过去位置不能看未来位置；模型参数在推理中也不更新。因此，对于历史位置 $i\le t$：

\[
\vect{k}_i^{(\ell)},\quad
\vect{v}_i^{(\ell)}
\]

在后续 decode 步中保持不变。

当生成第 $t+1$ 个 token 时，需要计算的是新 token 的 query：

\[
\vect{q}_{t+1}^{(\ell)}.
\]

它要和历史所有 key 做点积：

\[
\vect{q}*{t+1}^{(\ell)}
\left[
\vect{k}*{1}^{(\ell)},
\vect{k}*{2}^{(\ell)},
\ldots,
\vect{k}*{t+1}^{(\ell)}
\right]^{\mathsf{T}}.
\]

历史 key/value 不变，因此可以缓存。若不缓存，就必须在每个 decode step 重新对整个上下文做 forward，重复计算所有历史 token 的 K/V。这会把 decode 成本从每步近似 $O(t)$ 变成重复的 $O(t^2)$ 或更高，完全不可接受。

为什么不缓存 Query？因为每一步只需要当前新 token 的 query。历史 token 的 query 不会再用于未来 token 的 attention。未来 token 关注历史 token 时，需要的是未来 token 的 query 与历史 token 的 key，而不是历史 token 的 query。因此 Query 没有长期缓存价值。

可以总结为：

\[
\text{Key: 历史 token 的可匹配地址，需要被未来 query 查询。}
\]
\[
\text{Value: 历史 token 的内容，需要被未来 attention 加权汇聚。}
\]
\[
\text{Query: 当前 token 的查询向量，只在当前步使用。}
\]

\subsection{cache shape 推导}
\label{subsec:kv-cache-shape-derivation}

设模型有：

\begin{itemize}
\item 层数 $L$；
\item batch size $B$；
\item 当前最大上下文长度 $S$；
\item query head 数 $n_q$；
\item key/value head 数 $n_{kv}$；
\item head dimension $d_h$；
\item hidden size $H=n_qd_h$；
\item 每个元素字节数 $s_{\mathrm{dtype}}$。
\end{itemize}

对于每一层，K Cache 的 shape 通常是：

\[
\mat{K}*{cache}^{(\ell)}
\in
\mathbb{R}^{B\times n*{kv}\times S\times d_h}.
\]

V Cache 同理：

\[
\mat{V}*{cache}^{(\ell)}
\in
\mathbb{R}^{B\times n*{kv}\times S\times d_h}.
\]

所有层合并后：

\[
\mat{K}*{cache}
\in
\mathbb{R}^{L\times B\times n*{kv}\times S\times d_h},
\]
\[
\mat{V}*{cache}
\in
\mathbb{R}^{L\times B\times n*{kv}\times S\times d_h}.
\]

有些实现会使用不同 layout，例如：

\[
[B,L,n_{kv},S,d_h],
\]
\[
[L,B,S,n_{kv},d_h],
\]
\[
[L,B,n_{kv},d_h,S].
\]

layout 的选择会影响 attention kernel 的访存连续性、append 新 token 的写入效率、batch 内请求重排成本以及分页 KV Cache 的实现方式。

一个简单连续 KV Cache layout 如下图所示。

\begin{figure}[H]
\centering
\begin{tikzpicture}[
font=\small,
cell/.style={draw, minimum width=0.45cm, minimum height=0.38cm},
label/.style={font=\scriptsize},
box/.style={draw, rounded corners=3pt, minimum width=2.0cm, minimum height=0.65cm, align=center},
arrow/.style={-{Latex[length=2.1mm]}, thick}
]

\node[font=\bfseries] at (0,3.25) {KV Cache 张量布局示意};

\node[box, fill=blue!6] (layer) at (-4.7,1.9) {Layer $\ell$};
\node[box, fill=orange!12] (kv) at (-2.4,1.9) {K Cache / V Cache};
\draw[arrow] (layer) -- (kv);

\foreach \b/\y in {0/1.15,1/0.35,2/-0.45} {
  \node[label] at (-5.1,\y) {batch \b};
  \foreach \h/\xbase in {0/-3.95,1/-2.85,2/-1.75} {
    \node[label] at (\xbase,1.65) {head \h};
    \foreach \s in {0,...,4} {
      \node[cell, fill=green!10] at ({\xbase+0.18*\s},{\y}) {};
    }
  }
}

\draw[decorate,decoration={brace,amplitude=5pt},thick]
  (-4.05,-0.95) -- (-3.15,-0.95)
  node[midway,below=6pt] {$S$ positions};

\draw[decorate,decoration={brace,amplitude=5pt},thick]
  (-4.3,1.35) -- (-1.25,1.35)
  node[midway,above=6pt] {$n_{kv}$ heads};

\node[box, fill=purple!10, minimum width=2.4cm] (shape) at (2.6,0.4)
  {$[L,B,n_{kv},S,d_h]$};

\draw[arrow] (-1.1,0.4) -- (shape);

\node[align=left, text width=10.5cm] at (0,-2.05) {
  KV Cache 需要为每一层保存 K 和 V。每个请求的 cache 会随着生成长度增长。
  实际系统会根据 attention kernel 的访存模式选择具体内存 layout。
};

\end{tikzpicture}
\caption{KV Cache 的基本张量布局}
\label{fig:kv-cache-tensor-layout}
\end{figure}

\subsection{batch、head、seq\_len、head\_dim}
\label{subsec:batch-head-seqlen-headdim}

KV Cache shape 中每个维度都有系统含义。

\textbf{Batch 维 $B$} 对应同一时刻正在服务的请求数。与训练不同，推理中的 batch 是动态的。请求可能随时加入或退出，且每个请求的上下文长度不同。因此 KV Cache 的 batch 维并不总是一个固定矩形张量那样简单。传统实现可能为每个 batch slot 预留最大长度；更高效的实现会使用 block-based 或 paged allocation。

\textbf{Head 维 $n_{kv}$} 对应 key/value heads 数量。对于 MHA：

\[
n_{kv}=n_q.
\]

对于 MQA：

\[
n_{kv}=1.
\]

对于 GQA：

\[
1<n_{kv}<n_q.
\]

KV Cache 显存与 $n_{kv}$ 成正比。因此 GQA/MQA 对推理显存和带宽非常重要。

\textbf{Sequence 维 $S$} 对应每个请求已经缓存的 token 数。Prefill 后，$S$ 等于 prompt length；decode 每生成一个 token，$S$ 增加 1。对于多轮对话，请求的 KV Cache 可能不断增长，直到达到上下文窗口上限或触发截断、压缩、eviction。

\textbf{Head dimension $d_h$} 是每个 attention head 的向量维度。通常：

\[
H=n_qd_h.
\]

KV Cache 显存与 $d_h$ 成正比。对固定 hidden size $H$，如果使用 GQA 减少 $n_{kv}$，KV Cache 可显著下降。

\subsection{KV Cache 显存估算}
\label{subsec:kv-cache-memory-estimation}

KV Cache 显存是 serving capacity 的核心约束。对于一个 dense Decoder-only 模型，KV Cache 总显存可粗略估算为：

\[
M_{\mathrm{KV}}
2
\cdot
L
\cdot
B
\cdot
S
\cdot
n_{kv}
\cdot
d_h
\cdot
s_{\mathrm{dtype}}.
\]

其中前面的 $2$ 表示 K 和 V 两份 cache。

如果使用 tensor parallel，KV Cache 也可能按 head 或 hidden 维切分到多个 GPU 上。设 TP size 为 $P_{\mathrm{TP}}$，如果 $n_{kv}$ 能按 TP 维切分，则每 GPU KV Cache 近似为：

\[
M_{\mathrm{KV,perGPU}}
\frac{
2
\cdot
L
\cdot
B
\cdot
S
\cdot
n_{kv}
\cdot
d_h
\cdot
s_{\mathrm{dtype}}
}
{P_{\mathrm{TP}}}.
\]

如果使用 MHA，$n_{kv}=n_q$；如果使用 GQA，$n_{kv}$ 更小。以 $L=80$、$B=32$、$S=4096$、$n_{kv}=8$、$d_h=128$、BF16 为例：

\[
M_{\mathrm{KV}}
2
\cdot
80
\cdot
32
\cdot
4096
\cdot
8
\cdot
128
\cdot
2
\ \mathrm{bytes}.
\]

这个值约为：

\[
42.95\ \mathrm{GB}.
\]

这只是 KV Cache，不包括模型参数、临时 buffer、CUDA runtime、scheduler metadata 和其他开销。若 batch 更大、上下文更长，KV Cache 会迅速成为主要显存占用。

下面给出一个 KV Cache 显存估算器。

\begin{lstlisting}[language=Python,caption={KV Cache 显存估算器},label={lst:kv-cache-memory-estimator}]
def estimate_kv_cache_gb(
num_layers,
batch_size,
seq_len,
num_kv_heads,
head_dim,
dtype_bytes=2,
tensor_parallel_size=1,
):
"""
Estimate per-GPU KV cache memory in GB.

KV cache shape:
  K: [L, B, n_kv, S, d_h]
  V: [L, B, n_kv, S, d_h]
"""
total_bytes = (
    2
    * num_layers
    * batch_size
    * seq_len
    * num_kv_heads
    * head_dim
    * dtype_bytes
)

per_gpu_bytes = total_bytes / tensor_parallel_size
return per_gpu_bytes / (1024 ** 3)

gb = estimate_kv_cache_gb(
num_layers=80,
batch_size=32,
seq_len=4096,
num_kv_heads=8,
head_dim=128,
dtype_bytes=2,
tensor_parallel_size=1,
)

print(round(gb, 2), "GB")
\end{lstlisting}

KV Cache 显存估算有几个常见误区：

\begin{itemize}
\item \textbf{只看模型参数，不看 KV Cache}：对于长上下文和高并发，KV Cache 可能比参数更限制 batch size；
\item \textbf{忽略 GQA/MQA}：$n_{kv}$ 直接影响 KV Cache；
\item \textbf{忽略 TP 切分}：多卡推理时每 GPU cache 可能被切分；
\item \textbf{忽略碎片}：连续分配可能造成大量未使用空间；
\item \textbf{忽略最大长度预留}：为最大上下文预留会比实际 token 使用更多显存；
\item \textbf{忽略多轮对话增长}：同一个会话可能随着轮次持续增长。
\end{itemize}

\section{KV Cache 的性能瓶颈}
\label{sec:kv-cache-performance-bottlenecks}

KV Cache 不只是显存问题，也是性能问题。Decode 阶段每生成一个 token，都要在每一层读取历史 K/V。随着上下文变长，读取的数据量线性增长。由于每一步只处理一个新 token，计算密度不高，decode 往往受 HBM 带宽、cache layout、batch 组织和 kernel 访存效率限制。

\subsection{decode 阶段的 memory-bound 特性}
\label{subsec:decode-memory-bound}

在 decode 阶段，每层 attention 对当前 token 执行：

\[
\vect{q}*t\mat{K}*{cache}^{\mathsf{T}},
\]
\[
\softmax(\cdot)\mat{V}_{cache}.
\]

对于每个新 token、每层、每个 KV head，需要读取长度为 $S$ 的 K 和 V：

\[
\mat{K}*{cache}\in\mathbb{R}^{S\times d_h},
\quad
\mat{V}*{cache}\in\mathbb{R}^{S\times d_h}.
\]

读取量随 $S$ 线性增长：

\[
O(Sd_h).
\]

而当前 query 只有一个 token，矩阵乘法形状通常较小。相对于训练或 prefill 中的大 GEMM，decode attention 的算术强度较低：

\[
\text{Arithmetic Intensity}
\frac{\text{FLOPs}}
{\text{Bytes read/written}}.
\]

当算术强度低时，性能更受内存带宽限制，而不是 Tensor Core 峰值 FLOPs 限制。也就是说，即使 GPU 有很高的理论算力，decode 阶段也可能因为不断读取 KV Cache 而无法充分利用计算单元。

\begin{figure}[H]
\centering
\begin{tikzpicture}[
font=\small,
box/.style={draw, rounded corners=3pt, minimum width=1.65cm, minimum height=0.65cm, align=center},
mem/.style={draw, rounded corners=4pt, minimum width=2.2cm, minimum height=1.0cm, align=center, fill=orange!14},
arrow/.style={-{Latex[length=2.1mm]}, thick}
]

\node[font=\bfseries] at (0,3.25) {Decode 阶段为何容易 memory-bound};

\node[box, fill=blue!6] (q) at (-4.8,1.2) {new token\\$q_t$};
\node[mem] (kvcache) at (-1.8,1.2) {KV Cache\\history $1:t$};
\node[box, fill=green!10] (attn) at (1.3,1.2) {Attention\\small compute};
\node[box, fill=purple!10] (out) at (4.2,1.2) {next hidden};

\draw[arrow] (q) -- (attn);
\draw[arrow] (kvcache) -- node[above] {large HBM reads} (attn);
\draw[arrow] (attn) -- (out);

\node[box, fill=yellow!16] (ffn) at (1.3,-0.1) {FFN / Linear\\small batch GEMM};
\draw[arrow] (out) -- (ffn);

\node[align=left, text width=10.5cm] at (0,-1.65) {
  Decode 每步只处理新 token，计算规模小；但必须读取历史 K/V。
  上下文越长，KV Cache 读取越大，attention 越容易受 HBM 带宽限制。
};

\end{tikzpicture}
\caption{Decode 阶段的 memory-bound 特性}
\label{fig:decode-memory-bound}
\end{figure}

\subsection{长上下文显存压力}
\label{subsec:long-context-memory-pressure}

KV Cache 显存与上下文长度 $S$ 线性相关：

\[
M_{\mathrm{KV}}\propto S.
\]

如果上下文长度从 $4096$ 增加到 $32768$，KV Cache 显存也增加 $8$ 倍。长上下文模型不仅 prefill 更贵，decode 时每步读取的 cache 也更大。

设单请求 KV Cache 显存为：

\[
M_{\mathrm{KV,req}}
2L S n_{kv}d_hs_{\mathrm{dtype}}.
\]

在固定 GPU 显存预算 $M_{\mathrm{budget}}$ 下，可并发请求数近似满足：

\[
B_{\max}
\le
\frac{M_{\mathrm{budget}}}
{M_{\mathrm{KV,req}}}.
\]

这说明长上下文会直接降低可承载并发。对于在线服务，用户请求长度差异很大：有些请求只有几十 token，有些请求有几万 token。如果简单按最大长度预留 KV Cache，会浪费大量显存；如果按实际长度动态分配，又需要复杂内存管理。

长上下文还会影响调度。一个很长的请求可能占用大量 KV Cache，降低 batch 中其他请求的可用空间。服务系统必须决定是否接收、排队、截断、分块处理或拒绝该请求。这属于 admission control 的问题，后续 continuous batching 章节会继续讨论。

\subsection{batch size 与 cache 碎片}
\label{subsec:batch-size-cache-fragmentation}

推理中 batch size 不像训练那样固定。请求动态加入和退出，每个请求的 prompt length 和 output length 都不同。若使用简单连续分配，为每个请求预留最大长度，会产生两类浪费：

\begin{itemize}
\item \textbf{内部碎片}：请求实际使用长度小于预留长度，预留空间未使用；
\item \textbf{外部碎片}：已释放的空间分散，难以满足新的大请求连续分配。
\end{itemize}

例如两个请求：

\[
R_1:\ \text{最大预留 }4096,\quad \text{实际使用 }512,
\]
\[
R_2:\ \text{最大预留 }4096,\quad \text{实际使用 }3800.
\]

如果都按 $4096$ 预留，$R_1$ 浪费大量空间。如果按实际长度增长式分配，则当 $R_2$ 继续生成时，可能需要扩展连续空间；若后面空间已被其他请求占用，就需要搬移或重新分配。

这正是 PagedAttention 和 block-based allocation 的动机。把 KV Cache 切成固定大小 block，可以像操作系统分页一样管理，避免要求每个请求拥有连续大块空间。本章先讨论基本策略，下一章会深入 PagedAttention。

\subsection{多轮对话中的 cache 复用}
\label{subsec:multi-turn-cache-reuse}

多轮对话中，用户和助手的历史消息会不断追加。例如：

\[
\text{System prompt}
+
\text{User}_1
+
\text{Assistant}_1
+
\text{User}_2
+
\cdots
\]

如果每一轮都从头 prefill 整个对话历史，会重复计算大量相同前缀。KV Cache 可以在同一会话内复用历史前缀，只对新增用户消息做增量 prefill，再进入 decode。

假设前一轮已经缓存了前缀长度 $S_{\mathrm{old}}$：

\[
x_1,\ldots,x_{S_{\mathrm{old}}}.
\]

新一轮用户追加了 $\Delta S$ 个 token：

\[
x_{S_{\mathrm{old}}+1},\ldots,x_{S_{\mathrm{old}}+\Delta S}.
\]

系统只需要对新增部分做 prefill，并让新增 token attend 到旧 cache：

\[
\mat{K}*{cache}[1:S*{\mathrm{old}}+\Delta S].
\]

这样可以降低多轮对话中的重复 prefill 成本。但 cache 复用也有工程问题：

\begin{itemize}
\item 会话 cache 需要在请求之间保留，占用显存；
\item 用户可能长时间不继续对话，cache 是否保留需要 eviction 策略；
\item 多副本部署时，后续请求是否路由到同一 worker 影响复用；
\item 如果 system prompt 或采样参数变化，prefix 是否仍可复用需要判断；
\item 对话历史超出上下文窗口时，需要截断或压缩。
\end{itemize}

\begin{figure}[H]
\centering
\begin{tikzpicture}[
font=\small,
req/.style={draw, rounded corners=3pt, minimum width=1.8cm, minimum height=0.55cm, align=center},
cache/.style={draw, rounded corners=2pt, minimum width=0.65cm, minimum height=0.4cm, align=center},
arrow/.style={-{Latex[length=2.1mm]}, thick}
]

\node[font=\bfseries] at (0,3.25) {多请求 KV Cache 随生成过程增长};

\foreach \r/\y/\len/\col in {1/1.65/4/blue!10,2/0.75/7/green!10,3/-0.15/3/orange!14} {
  \node[req, fill=\col] at (-5.0,\y) {Request \r};
  \foreach \i in {1,...,\len} {
    \node[cache, fill=\col] at ({-3.5+0.45*\i},\y) {};
  }
  \draw[arrow] ({-3.5+0.45*\len+0.35},\y) -- ({-3.5+0.45*\len+1.0},\y);
  \node[cache, fill=yellow!18] at ({-3.5+0.45*\len+1.3},\y) {new};
}

\node[align=left, text width=10.5cm] at (0,-1.45) {
  不同请求的 KV Cache 长度不同，并且 decode 每步继续增长。
  Serving 系统需要动态管理 cache 空间，避免长短请求混部造成显存碎片和调度低效。
};

\end{tikzpicture}
\caption{多请求 KV Cache 的增长过程}
\label{fig:multi-request-kv-cache-growth}
\end{figure}

\section{KV Cache 管理策略}
\label{sec:kv-cache-management-strategies}

KV Cache 管理是 LLM serving 系统的核心模块之一。它不仅负责分配和释放显存，还要配合 scheduler、attention kernel、prefix cache、请求取消、长上下文、batch 重排和多 GPU 并行。

本节讨论四类基础策略：

\begin{itemize}
\item contiguous allocation；
\item block-based allocation；
\item cache eviction；
\item prefix caching。
\end{itemize}

\subsection{contiguous allocation}
\label{subsec:contiguous-allocation}

最简单的 KV Cache 管理方式是连续分配。系统为每个请求分配一段连续的 cache 空间：

\[
[\text{start},\ \text{start}+S_{\max}).
\]

每生成一个 token，就把新的 K/V 写入下一个位置。连续分配的优点是实现简单，attention kernel 读取也简单，因为一个请求的历史 K/V 在内存中是连续的。

连续分配的常见布局可以是：

\[
\texttt{cache[layer][batch_slot][kv_head][pos][head_dim]}.
\]

写入第 $t$ 个 token 时：

\[
\texttt{cache[layer][slot][head][t][:]}
\leftarrow
\texttt{new_kv}.
\]

这种方式适合固定 batch、固定最大长度、请求长度较一致的场景。但在线服务中，连续分配有明显缺点：

\begin{itemize}
\item 需要提前预留最大长度，浪费显存；
\item 请求长度不均导致内部碎片；
\item 请求释放后产生外部碎片；
\item 扩展请求长度可能需要重新分配；
\item batch slot 与 request 绑定，动态调度不灵活。
\end{itemize}

\begin{figure}[H]
\centering
\begin{tikzpicture}[
font=\small,
cell/.style={draw, minimum width=0.48cm, minimum height=0.42cm},
label/.style={font=\scriptsize}
]

\node[font=\bfseries] at (0,3.0) {Contiguous KV Cache Allocation};

\foreach \r/\y/\used/\col in {1/1.55/3/blue!10,2/0.85/6/green!10,3/0.15/2/orange!14} {
  \node[label] at (-4.8,\y) {Request \r};
  \foreach \i in {0,...,7} {
    \ifnum\i<\used
      \node[cell, fill=\col] at ({-3.8+0.52*\i},\y) {};
    \else
      \node[cell, fill=gray!16] at ({-3.8+0.52*\i},\y) {};
    \fi
  }
}

\node[draw, rounded corners=2pt, fill=blue!10, minimum width=0.55cm, minimum height=0.35cm] at (1.4,1.5) {};
\node[anchor=west] at (1.8,1.5) {used tokens};
\node[draw, rounded corners=2pt, fill=gray!16, minimum width=0.55cm, minimum height=0.35cm] at (1.4,0.9) {};
\node[anchor=west] at (1.8,0.9) {reserved but unused};

\node[align=left, text width=10.4cm] at (0,-1.35) {
  连续分配简单、访存友好，但需要预留最大长度。
  在线请求长度不均时，未使用的预留空间会造成显著显存浪费。
};

\end{tikzpicture}
\caption{连续 KV Cache 分配中的预留空间浪费}
\label{fig:contiguous-kv-cache-allocation}
\end{figure}

\subsection{block-based allocation}
\label{subsec:block-based-allocation}

Block-based allocation 将 KV Cache 切成固定大小的 blocks。每个 block 可以保存固定数量 token 的 K/V，例如 block size 为 $16$ 或 $32$ tokens。请求不再需要一段连续大空间，而是持有一个 block 列表：

\[
\text{Request }r
\rightarrow
[b_0,b_1,b_2,\ldots].
\]

当请求继续生成，需要更多空间时，从 free block pool 中申请新 block；请求结束时，把它持有的 blocks 释放回池中。

这种方式的优点是：

\begin{itemize}
\item 避免按最大长度连续预留；
\item 释放后 block 可被其他请求复用；
\item 更适合动态 batch 和长短请求混部；
\item 支持 prefix cache 的 block 级共享；
\item 为 PagedAttention 的 block table 奠定基础。
\end{itemize}

缺点是 attention kernel 读取 K/V 时不能简单假设请求的历史 token 连续存储。Kernel 需要通过 block table 把逻辑位置映射到物理 block。这会增加地址计算复杂度，但换来显存利用率提升。

\begin{figure}[H]
\centering
\begin{tikzpicture}[
font=\small,
block/.style={draw, rounded corners=2pt, minimum width=0.7cm, minimum height=0.45cm, align=center},
req/.style={draw, rounded corners=3pt, minimum width=1.55cm, minimum height=0.55cm, align=center, fill=blue!6},
arrow/.style={-{Latex[length=2.1mm]}, thick}
]

\node[font=\bfseries] at (0,3.15) {Block-based KV Cache Allocation};

\node[req] (r1) at (-5.2,1.55) {Req 1};
\node[req] (r2) at (-5.2,0.55) {Req 2};
\node[req] (r3) at (-5.2,-0.45) {Req 3};

\foreach \i/\x/\y/\col in {
  0/-2.8/1.55/blue!10,
  1/-1.85/1.55/blue!10,
  2/-0.9/1.55/blue!10,
  3/-2.8/0.55/green!10,
  4/-1.85/0.55/green!10,
  5/-0.9/0.55/green!10,
  6/0.05/0.55/green!10,
  7/-2.8/-0.45/orange!14,
  8/-1.85/-0.45/orange!14
} {
  \node[block, fill=\col] (b\i) at (\x,\y) {$b_{\i}$};
}

\draw[arrow] (r1) -- (b0);
\draw[arrow] (r2) -- (b3);
\draw[arrow] (r3) -- (b7);

\node[draw, rounded corners=3pt, fill=gray!12, minimum width=2.9cm, minimum height=0.75cm, align=center] (free) at (3.1,0.55)
  {Free block pool};

\foreach \i/\x in {9/2.25,10/3.1,11/3.95} {
  \node[block, fill=gray!20] at (\x,-0.25) {$b_{\i}$};
}

\node[align=left, text width=10.5cm] at (0,-1.95) {
  每个请求持有若干固定大小 block，而不是一段连续大空间。
  新 token 写满当前 block 后，再从 free pool 申请新 block。
};

\end{tikzpicture}
\caption{Block-based KV Cache 分配}
\label{fig:block-based-kv-cache-allocation}
\end{figure}

\subsection{cache eviction}
\label{subsec:cache-eviction}

KV Cache 占用显存，而显存是有限资源。当请求结束、取消、超时或被调度策略移除时，需要释放 cache。更复杂的是，系统可能为了 prefix caching 或多轮会话保留某些 cache，以便未来复用。此时需要 eviction 策略决定哪些 cache 可以被驱逐。

常见 eviction 依据包括：

\begin{itemize}
\item \textbf{请求是否已结束}：结束请求的 cache 可立即释放；
\item \textbf{会话是否仍活跃}：长时间无后续请求的 session cache 可释放；
\item \textbf{prefix 复用频率}：高频 system prompt 或热门 prefix 值得保留；
\item \textbf{cache 大小}：超大 cache 可能优先驱逐以释放显存；
\item \textbf{SLA 优先级}：低优先级请求的 cache 可更早释放；
\item \textbf{最近使用时间}：LRU 类策略；
\item \textbf{预计重算成本}：prefill 很贵的 prefix 更值得缓存。
\end{itemize}

可以定义一个 cache value 分数：

\[
\operatorname{score}(c)
\alpha\cdot \operatorname{reuse_prob}(c)
+
\beta\cdot \operatorname{recompute_cost}(c)
-------------------------------------------
## \gamma\cdot \operatorname{memory_size}(c)
\delta\cdot \operatorname{idle_time}(c).
\]

分数低的 cache 更适合驱逐。真实系统不一定显式使用这样的公式，但背后的权衡类似：保留 cache 可以降低未来 prefill 成本，但会占用当前显存，降低当前可服务并发。

\subsection{prefix caching}
\label{subsec:prefix-caching}

Prefix caching 指多个请求共享相同前缀时，复用该前缀的 KV Cache。典型场景包括：

\begin{itemize}
\item 所有请求共享同一个 system prompt；
\item RAG 模板中前半部分相同；
\item 多个用户请求使用相同 few-shot examples；
\item 多轮对话中复用历史上下文；
\item 批量评测中 prompt 模板一致。
\end{itemize}

假设两个请求前缀相同：

\[
x_1,\ldots,x_K.
\]

它们的 KV Cache 前 $K$ 个 token 完全相同。因此可以只保存一份 prefix KV，并让多个请求引用它。每个请求只为自己的 suffix 分配新的 cache：

\[
\text{Request A}
\text{Shared Prefix}
+
\text{Suffix A},
\]
\[
\text{Request B}
\text{Shared Prefix}
+
\text{Suffix B}.
\]

Prefix cache 的收益是减少重复 prefill 和重复显存占用。但它带来引用计数和生命周期管理问题。共享 prefix 不能在仍有请求引用时释放；如果某个请求继续生成，只能追加自己的 suffix，不能修改共享 prefix。

\begin{figure}[H]
\centering
\begin{tikzpicture}[
font=\small,
block/.style={draw, rounded corners=2pt, minimum width=0.6cm, minimum height=0.42cm, align=center},
req/.style={draw, rounded corners=3pt, minimum width=1.45cm, minimum height=0.55cm, align=center, fill=blue!6},
arrow/.style={-{Latex[length=2.1mm]}, thick}
]

\node[font=\bfseries] at (0,3.25) {Prefix Cache 复用示意};

\node[req] (ra) at (-5.2,1.25) {Req A};
\node[req] (rb) at (-5.2,0.15) {Req B};

\foreach \i/\x in {0/-2.9,1/-2.25,2/-1.6,3/-0.95} {
  \node[block, fill=green!12] (p\i) at (\x,1.85) {$p_{\i}$};
}
\draw[decorate,decoration={brace,amplitude=5pt},thick]
  (-3.25,2.25) -- (-0.6,2.25)
  node[midway,above=6pt] {shared prefix KV};

\foreach \i/\x in {0/0.4,1/1.05,2/1.7} {
  \node[block, fill=blue!10] (a\i) at (\x,1.25) {$a_{\i}$};
}
\foreach \i/\x in {0/0.4,1/1.05} {
  \node[block, fill=orange!14] (b\i) at (\x,0.15) {$b_{\i}$};
}

\draw[arrow] (ra) -- (p0);
\draw[arrow] (rb) -- (p0);

\draw[arrow] (p3) -- (a0);
\draw[arrow] (p3) -- (b0);

\node[align=left, text width=10.5cm] at (0,-1.25) {
  多个请求可以共享相同 prefix 的 KV Cache，只为各自 suffix 分配额外 cache。
  这能降低重复 prefill 和显存占用，但需要引用计数与生命周期管理。
};

\end{tikzpicture}
\caption{Prefix cache 在多个请求之间复用相同前缀}
\label{fig:prefix-cache-reuse}
\end{figure}

下面给出一个极简 prefix cache 管理伪代码。

\begin{lstlisting}[language=Python,caption={Prefix Cache 管理的极简伪代码},label={lst:prefix-cache-manager}]
class PrefixCacheEntry:
def **init**(self, blocks):
self.blocks = blocks
self.ref_count = 0
self.last_used_time = 0

class PrefixCacheManager:
def **init**(self):
self.table = {}  # prefix_hash -> PrefixCacheEntry

def lookup(self, prefix_tokens):
    key = hash_tokens(prefix_tokens)
    entry = self.table.get(key)
    if entry is None:
        return None

    entry.ref_count += 1
    entry.last_used_time = now()
    return entry.blocks

def insert(self, prefix_tokens, blocks):
    key = hash_tokens(prefix_tokens)
    entry = PrefixCacheEntry(blocks)
    entry.ref_count = 1
    entry.last_used_time = now()
    self.table[key] = entry
    return entry.blocks

def release(self, prefix_tokens):
    key = hash_tokens(prefix_tokens)
    entry = self.table.get(key)
    if entry is not None:
        entry.ref_count -= 1

def evict_if_needed(self, free_block_pool):
    # Evict unused prefixes first.
    candidates = [
        (key, entry)
        for key, entry in self.table.items()
        if entry.ref_count == 0
    \]

    # Example policy: least recently used.
    candidates.sort(key=lambda item: item[1].last_used_time)

    for key, entry in candidates:
        if free_block_pool.has_enough_free_blocks():
            break
        free_blocks(entry.blocks)
        del self.table[key]

\end{lstlisting}

真实 prefix caching 要处理更多细节，例如 tokenizer 一致性、prefix hash 冲突、block table 共享、copy-on-write、并发请求、跨 worker cache 命中、cache warmup 和安全隔离。

\section{KV Cache 与并行推理}
\label{sec:kv-cache-and-parallel-inference}

前面主要从单 GPU 或单模型副本角度讨论 KV Cache。实际 serving 中，大模型常使用 tensor parallel、pipeline parallel 或 expert parallel 进行多 GPU 推理。此时 KV Cache 也必须与并行方式匹配。

\subsection{Tensor Parallel 下的 KV Cache}
\label{subsec:kv-cache-under-tensor-parallel}

Tensor Parallel 推理中，attention heads 通常按 TP rank 切分。若 query heads 和 kv heads 都能整除 TP size，则每个 rank 保存一部分 KV heads：

\[
n_{kv,\mathrm{local}}
\frac{n_{kv}}{P_{\mathrm{TP}}}.
\]

每 GPU KV Cache 近似为：

\[
M_{\mathrm{KV,local}}
2
\cdot
L
\cdot
B
\cdot
S
\cdot
n_{kv,\mathrm{local}}
\cdot
d_h
\cdot
s_{\mathrm{dtype}}.
\]

也就是：

\[
M_{\mathrm{KV,local}}
\frac{M_{\mathrm{KV,total}}}{P_{\mathrm{TP}}}.
\]

在这种设计中，每个 rank 本地执行自己负责 heads 的 attention，读取本地 KV Cache。之后 attention output projection 通常需要 TP all-reduce 或其他通信恢复 hidden state。

如果使用 GQA，需要保证 kv heads 与 TP 的映射合理。若 $n_{kv}$ 小于 TP size，某些 TP rank 可能没有独立 kv head，或者需要复制 K/V。实现上可能有多种策略：

\begin{itemize}
\item 让 TP size 不超过 $n_{kv}$；
\item 在多个 query ranks 之间共享 kv heads；
\item 复制 K/V 到多个 ranks；
\item 使用更复杂的 attention kernel mapping。
\end{itemize}

这说明模型架构中的 GQA 参数会影响 serving 并行配置。

\subsection{Pipeline Parallel 下的 KV Cache}
\label{subsec:kv-cache-under-pipeline-parallel}

Pipeline Parallel 推理中，不同 GPU 负责不同层。每个 pipeline stage 只保存自己负责层的 KV Cache。若第 $p$ 个 stage 负责层集合 $\mathcal{L}_p$，则它保存：

\[
{\mat{K}^{(\ell)},\mat{V}^{(\ell)}\mid \ell\in\mathcal{L}_p}.
\]

每个 token 在 decode 时依次经过所有 stages。Stage 0 处理完后把 hidden state 发送给 Stage 1，直到最后 stage 输出 logits。

PP 推理中的 KV Cache 有几个问题：

\begin{itemize}
\item 每个 stage 的 cache 显存取决于它负责的层数；
\item stage 间需要传递每个 decode step 的 hidden state；
\item decode token-by-token 使 pipeline bubble 更明显；
\item 如果 batch 很小，PP 利用率可能较低；
\item pipeline stage 负载不均会影响 TPOT。
\end{itemize}

训练中的 pipeline 依靠多个 micro-batches 填充流水线；推理 decode 中，每个请求每步只有一个 token，流水线填充更加困难。因此很多推理系统更偏向 tensor parallel，而不是过深 pipeline parallel。对于非常大的模型，PP 仍可能必要，但需要结合 continuous batching 提高利用率。

\subsection{KV Cache 与多副本部署}
\label{subsec:kv-cache-and-replica-serving}

在线服务通常部署多个模型副本。一个用户会话的 KV Cache 常驻在某个 worker 上。如果下一轮请求被路由到另一个 worker，就无法直接复用该 cache，除非做跨 worker cache 迁移或共享存储。

这带来 sticky session 问题：

\[
\text{same session}
\rightarrow
\text{same worker}
\]

可以提高 cache 命中率，但可能导致负载不均。例如某些 worker 上积累了很多长会话，显存压力更高；另一些 worker 空闲。若为了负载均衡把会话迁移到其他 worker，则可能失去 KV Cache 复用，增加 prefill 成本。

因此 serving router 需要在以下目标之间权衡：

\begin{itemize}
\item cache locality；
\item worker load balance；
\item GPU memory pressure；
\item request priority；
\item tail latency；
\item session affinity。
\end{itemize}

\section{KV Cache 实现骨架}
\label{sec:kv-cache-implementation-skeleton}

下面给出一个教学版 KV Cache 类。它采用连续分配，不支持分页、prefix sharing 或动态 batch 退出，但有助于理解 cache 的基本操作：初始化、append、读取和长度维护。

\begin{lstlisting}[language=Python,caption={教学版连续 KV Cache 实现骨架},label={lst:kv-cache-class-skeleton}]
import torch

class ContiguousKVCache:
def **init**(
self,
num_layers,
batch_size,
max_seq_len,
num_kv_heads,
head_dim,
dtype=torch.bfloat16,
device="cuda",
):
self.num_layers = num_layers
self.batch_size = batch_size
self.max_seq_len = max_seq_len
self.num_kv_heads = num_kv_heads
self.head_dim = head_dim

    # K/V: [L, B, H_kv, S_max, D]
    self.k_cache = torch.empty(
        num_layers,
        batch_size,
        num_kv_heads,
        max_seq_len,
        head_dim,
        dtype=dtype,
        device=device,
    )
    self.v_cache = torch.empty_like(self.k_cache)

    # Current valid length for each batch slot.
    self.lengths = torch.zeros(
        batch_size,
        dtype=torch.long,
        device=device,
    )

def append(self, layer_idx, batch_indices, k_new, v_new):
    """
    Append K/V for one decode step.

    k_new, v_new: [B_active, H_kv, 1, D]
    batch_indices: [B_active]
    """
    for local_i, batch_i in enumerate(batch_indices):
        pos = self.lengths[batch_i].item()

        if pos >= self.max_seq_len:
            raise RuntimeError("KV cache is full")

        self.k_cache[layer_idx, batch_i, :, pos:pos+1, :] = k_new[local_i]
        self.v_cache[layer_idx, batch_i, :, pos:pos+1, :] = v_new[local_i]

def increment_lengths(self, batch_indices):
    self.lengths[batch_indices] += 1

def get_prefix(self, layer_idx, batch_i):
    """
    Return valid K/V prefix for one request.

    k: [H_kv, S, D]
    v: [H_kv, S, D]
    """
    length = self.lengths[batch_i].item()
    k = self.k_cache[layer_idx, batch_i, :, :length, :]
    v = self.v_cache[layer_idx, batch_i, :, :length, :]
    return k, v

def reset_slot(self, batch_i):
    # We do not need to clear memory immediately.
    # Updating length is enough for correctness.
    self.lengths[batch_i] = 0

\end{lstlisting}

这个实现有许多生产系统不能接受的问题：

\begin{itemize}
\item Python loop 不能用于高性能 decode；
\item batch slot 固定，动态加入退出不灵活；
\item 每个请求预留 \texttt{max\_seq\_len}，浪费显存；
\item 不支持 block table；
\item 不支持 prefix cache；
\item 不支持 paged attention kernel；
\item 不支持 tensor parallel 中的 head 切分；
\item 不支持请求取消后的 block 复用；
\item 不支持异步调度与 stream 管理。
\end{itemize}

但它清楚地展示了 KV Cache 的基本语义：每一层、每个 batch slot、每个 KV head、每个 position 都保存一条 K/V 向量。

生产级实现通常会将 append 和 attention 读取融合到 C++/CUDA/Triton kernel 中，并使用 block table 或 paged memory 避免连续分配问题。

\section{本章小结}
\label{sec:chapter12-summary}

本章系统介绍了 KV Cache 的机制、推理流程、显存估算、性能瓶颈和管理策略。

\begin{itemize}
\item \textbf{LLM 推理分为 prefill 和 decode 两个阶段}。Prefill 一次处理完整 prompt，建立 KV Cache；decode 每次生成一个 token，并持续读取和追加 KV Cache。
\item \textbf{Prefill 更像训练 forward}，计算密度较高，但长 prompt 会带来 $S^2$ attention 成本，并决定 TTFT。
\item \textbf{Decode 是 token-by-token 的强依赖过程}，每步只处理新 token，但需要读取历史所有 K/V，因此容易 memory-bound。
\item \textbf{Key 和 Value 可以缓存}，因为历史 token 在 causal decoder 中不会被未来 token 改变；未来 token 只需要用当前 query 查询历史 key，并加权历史 value。
\item \textbf{Query 通常不缓存}，因为历史 query 不会被未来 token 复用。
\item \textbf{KV Cache shape 通常为 $[L,B,n_{kv},S,d_h]$}，不同实现会根据 attention kernel 和内存访问模式选择不同 layout。
\item \textbf{KV Cache 显存公式} 为 $M_{\mathrm{KV}}=2LBSn_{kv}d_hs_{\mathrm{dtype}}$，其中 $2$ 表示 K 和 V。
\item \textbf{GQA/MQA 能显著降低 KV Cache}，因为 cache 大小与 $n_{kv}$ 成正比。
\item \textbf{Decode 阶段的主要瓶颈常是 HBM 带宽}，因为每步需要读取随上下文长度线性增长的 K/V。
\item \textbf{长上下文会线性增加 KV Cache 显存和 decode 读取量}，直接限制并发能力和成本。
\item \textbf{连续 KV Cache 分配简单但容易浪费显存}，尤其在请求长度不均和动态加入退出的在线 serving 场景中。
\item \textbf{Block-based allocation 更适合动态请求}，它把 cache 切成固定大小 blocks，为 PagedAttention 奠定基础。
\item \textbf{Cache eviction 与 prefix caching 是 serving 系统中的关键策略}，它们在显存占用、重算成本和请求延迟之间做权衡。
\item \textbf{并行推理中的 KV Cache 必须与 TP/PP 配合}：TP 常按 KV heads 切分 cache，PP 则按层切分 cache。
\end{itemize}

下一章将深入 PagedAttention 与 vLLM。我们会看到，KV Cache 最大的工程问题之一不是公式中的显存大小，而是在线请求中长度不均、动态增长和释放带来的碎片。PagedAttention 通过类似操作系统分页的 block table，把逻辑 token 位置映射到物理 cache blocks，从而显著提升显存利用率，并支撑 continuous batching。
